% VERSION MAP: v0.5 method-first (09-01) | v0.6 question-first + genealogy (09-03 Fable) | v0.61 narrative-first (09-03 Sol) | v0.7 = Sol spine + Fable body, PI rulings (09-03) | v0.71 = Sol red-team R2 adjudicated: WSI-CL genealogy, boundary, wording (09-03)
\documentclass{article}
\usepackage{spconf,amsmath,amssymb,graphicx,hyperref}
\usepackage{booktabs}
\usepackage{multirow}

\def\x{{\mathbf x}}
\def\L{{\cal L}}
\newcommand{\Fg}{F_g}
\newcommand{\Fp}{F_p}
\newcommand{\Psel}{\mathcal{P}}
\newcommand{\Mem}{\mathcal{M}}

\title{REMEMBERING WHERE TO LOOK: CONTINUAL HIERARCHICAL EVIDENCE SELECTION FOR WHOLE-SLIDE IMAGES}

\name{Hsi-Ren Hung \qquad [Second Author] \qquad Huei-Fang Yang\sthanks{corresponding author. This work was supported in part by the National Science and Technology Council, Taiwan under Grant NSTC 112-2221-E-110-047-MY3, 114-2640-E-110-004, and 114-2634-F-006-002.}}
\address{Department of Computer Science and Engineering, National Sun Yat-sen University, Taiwan}

\begin{document}
\ninept
\maketitle

% -----------------------------------------------------------------------------
\begin{abstract}
Whole-slide image (WSI) diagnosis is an evidence-acquisition problem: a model must select a few patches from thousands, and its prediction is only as reliable as the evidence it chooses.
Continual-learning (CL) methods usually preserve classifiers or representations, but a budgeted diagnostic system can fail in a different way: its diagnostic pathway may remain intact while its selector gradually forgets \emph{where to look}.
We introduce \textbf{PathSelect}, which makes the evidence-selection policy itself the object of continual learning.
To isolate this failure, the image encoder, text encoder, and cosine diagnostic head are frozen; only a hierarchical group-to-patch selector is learned.
Within a task, three supervision signals answer complementary questions: whether the selected evidence predicts correctly, whether it is semantically discriminative, and which candidate patch would counterfactually reduce diagnostic loss.
Across tasks, per-task LoRA residuals are merged exactly into one shared selector, while a feature-free Selection Memory preserves past behavior through dual-level distillation, a utility-preservation hinge, and replay.
The key distinction is between \emph{behavioral continuity}---how the selector scores evidence---and \emph{functional continuity}---whether newly selected evidence remains diagnostically useful.
On four sequential TCGA tasks, PathSelect raises final class-incremental accuracy by 34.6 points over sequential fine-tuning in every seed, halves forgetting relative to replay alone ($0.0539$ vs.\ $0.1102$), and retains two orders of magnitude more of its originally selected evidence, while inference uses one task-free shared selector.
\end{abstract}

\begin{keywords}
whole-slide images, continual learning, evidence selection, computational pathology, LoRA
\end{keywords}

% -----------------------------------------------------------------------------
\section{Introduction}
\label{sec:intro}

A gigapixel WSI can contain tens of thousands of patches, yet a diagnosis often depends on a small subset.
Practical systems therefore operate under an \emph{evidence budget}: choose $B$ patches, then decide.
This creates a learning problem that is easy to miss in conventional continual learning.
A model may preserve its classifier and representation reasonably well, yet after learning new disease cohorts it can still fail because it no longer chooses the tissue that once supported the correct decision.

We therefore ask a narrower question:
\begin{quote}
\emph{When diagnostic tasks arrive sequentially, does a model forget where to look, and can that ability be preserved, and its loss measured, directly?}
\end{quote}
PathSelect is designed around this question rather than around a collection of modules.
Its first principle is \textbf{attribution by freezing}.
CONCH image/text encoders \cite{lu2024conch} and the cosine diagnostic rule are fixed throughout training.
There is no trainable classifier.
The only learned object is an evidence-selection policy.
For paired methods evaluated at the same checkpoint and class set the diagnostic pathway is identical, so between-method differences arise from selector updates; within-method class-incremental changes over time additionally include label-space expansion. Direct selection metrics provide an additional head-independent measurement.

Our hierarchical substrate follows the coarse-to-fine tissue-aware idea of HistoSelect \cite{huang2026histoselect}: first allocate attention among tissue groups, then select patches within each group.
We do \emph{not} claim this hierarchy itself as our primary novelty.
Instead, PathSelect starts where such selectors normally stop: after a selector has learned where to look, how can one keep that ability when tasks arrive sequentially?

This perspective leads to a deliberate separation between two kinds of learning.
Within a task, PathSelect must \emph{learn where to look now}.
Across tasks, it must \emph{keep looking well}.
The first is trained by diagnostic, semantic, and counterfactual-utility supervision.
The second is trained by three preservation questions: (i) do I still score old evidence similarly, (ii) even if I select differently, is the new evidence at least as useful, and (iii) can I still solve the old diagnostic problem?

\textbf{Contributions.}
\begin{itemize}
    \item We formulate continual WSI diagnosis with the \emph{evidence-selection policy}---rather than a classifier---as the learned, preserved, and measured object, isolated by a frozen diagnostic pathway; we have not found prior WSI continual-learning work that freezes the complete diagnostic pathway and makes a hard, budgeted evidence-selection policy the sole learned and directly measured object.
    \item We introduce a fixed-cost continual protocol in which per-task LoRA residuals are merged into one shared hierarchical selector, while a feature-free Selection Memory preserves both \emph{selection behavior} and \emph{selection usefulness}.
    \item We pair the preservation objectives with selection-level measurements: Jaccard for identity drift and utility retention for functional drift, and show experimentally that these two notions are related but not interchangeable.
\end{itemize}

% -----------------------------------------------------------------------------
\section{Related Work}
\label{sec:related}

\textbf{Evidence selection in WSIs.}
MIL methods such as ABMIL, CLAM, and TransMIL \cite{ilse2018abmil,lu2021clam,shao2021transmil} aggregate patch features with learned attention, but selection is entangled with a trained diagnostic model.
HistoSelect \cite{huang2026histoselect} makes selection explicitly hierarchical, using tissue-aware coarse-to-fine reasoning.
PathSelect inherits this group$\rightarrow$patch inductive bias but changes the scientific setting: the diagnostic pathway is frozen and the selector is learned under a sequential task stream. HistoSelect conditions selection on a per-question query for VLM reasoning; PathSelect instead fixes a task-free tissue vocabulary, because a policy that must persist across tasks needs a grouping that does not move.

\textbf{LoRA-based continual learning.}
Low-rank adaptation \cite{hu2022lora} provides parameter-efficient task updates.
Prior CL methods use low-rank or modular updates to preserve classifier/backbone knowledge or retrieve task-specific components \cite{liang2024inflora,wu2025sdlora,he2025cllora,chen2026flexem}.
PathSelect instead applies continual learning only to the evidence selector.
Importantly, LoRA merging is used as a \emph{deployment substrate}, not claimed as the preservation mechanism: our experiments show that merging alone does not prevent forgetting.

\textbf{Continual WSI learning.} ConSlide \cite{huang2023conslide} introduced rehearsal-based continual WSI analysis on a four-organ TCGA sequence; MICIL \cite{meseguer2024micil} combined MIL with class-incremental learning; QPMIL-VL \cite{gou2025qpmil} brought vision--language prototypes to incremental WSI classification on CONCH features; and, closest to us, Li et al.\ \cite{li2025akd} located forgetting in the attention layers of MIL models and proposed attention distillation with a pseudo-bag memory of informative patches. All of these retain a trainable MIL diagnostic model in which attention is part of the aggregator. PathSelect differs on four axes: the entire diagnostic pathway is frozen and the selector is the sole learned object; selection is a hard, budgeted set rather than soft attention weights; preservation targets the utility of the selected set in addition to score behavior, so identity may change without functional regression; and selection identity and utility are measured directly rather than inferred from accuracy.

% -----------------------------------------------------------------------------
\section{Method}
\label{sec:method}

\subsection{Study object: a frozen diagnostic pathway}
\label{ssec:setup}

Tasks $\tau=1,\ldots,T$ arrive sequentially.
For a slide, frozen CONCH features are
\[
\{x_i\}_{i=1}^{N}, \qquad x_i\in\mathbb{R}^{512}.
\]
The selector must choose an evidence set
\[
\Psel=\{x_{i_1},\ldots,x_{i_B}\}, \qquad |\Psel|=B=8.
\]
Only this evidence reaches the diagnostic rule; the budget $B{=}8$ is the operating point at which learned selection peaked over random selection in a preliminary sweep (Sec.~\ref{sec:exp}).
If selector scores for the chosen patches are $s_i$, the WSI representation is the $\ell_2$-normalized score-weighted mean $z=\sum_{i\in\Psel}\mathrm{softmax}(s)_i\,x_i/\|\cdot\|_2$.
Prediction is the cosine similarity between $z$ and frozen class-text embeddings $\{t_c\}$, followed by softmax.
Training is task-aware; inference is task-free over all seen classes.
No task identity, adapter bank, memory lookup, or trained diagnostic head is used at test time.

\begin{figure*}[t]
    \centering
    \IfFileExists{figures/fig1_architecture_v1_0.pdf}{
        \includegraphics[width=\textwidth]{figures/fig1_architecture_v1_0.pdf}
    }{
        \fbox{\parbox[c][2.0cm][c]{0.96\textwidth}{\centering
        \textbf{Architecture figure placeholder.}\\
        Frozen CONCH features $\rightarrow$ semantic tissue groups $\rightarrow$
        group selector $\Fg$ $\rightarrow$ whole-budget quotas $\rightarrow$
        patch selector $\Fp$ $\rightarrow$ $B{=}8$ evidence $\rightarrow$
        frozen cosine diagnosis;\\
        across tasks: task LoRA $\rightarrow$ exact merge into shared selector;
        Selection Memory $\rightarrow$ KD / utility-equivalence / replay.}}
    }
    \caption{\textbf{PathSelect separates learning to look from keeping looking well.}
    Within a task, the hierarchical selector learns a budgeted evidence policy under a frozen diagnostic pathway.
    Across tasks, only selector LoRA parameters are trained; task residuals are folded into one shared selector.
    A feature-free Selection Memory is used only during training to preserve behavior and function.}
    \label{fig:arch}
\end{figure*}

\subsection{Hierarchical budgeted selection}
\label{ssec:hier}

Following the tissue-aware coarse-to-fine principle of HistoSelect \cite{huang2026histoselect}, patches are assigned semantically rather than clustered.
Eight fixed tissue prompts are encoded once with the frozen text encoder.
Each patch is assigned by
\[
G_j=\left\{x_i:
j=\arg\max_k \cos(x_i,p_k)\right\},
\]
and each group prototype is
\[
g_j=\operatorname{mean}(x_i\in G_j).
\]

Two small two-layer MLPs ($\Fg$ and $\Fp$, identical in structure and independent in weights) implement the policy; their input is zero-padded to a common width so that every ablation arm shares one architecture and parameter count.
The group selector $\Fg$ produces group scores $r_j$.
A largest-remainder allocation converts $\mathrm{softmax}(r)$ into integer quotas
\[
b_j\in\mathbb{N}, \qquad \sum_j b_j=B.
\]
The patch selector $\Fp$ produces patch scores $s_i$, and the top $b_j$ patches are selected within each group.
Thus the hierarchy answers two different questions:
\emph{which tissue groups deserve the budget?} and \emph{which patches within those groups deserve the slots?}

\subsection{Learning where to look: three within-task supervision signals}
\label{ssec:within}

The within-task objective is
\begin{equation}
\L_{\mathrm{ev}}
=
\L_{\mathrm{diag}}
+\beta_s\L_{\mathrm{sem}}
+\beta_u\L_{\mathrm{util}}.
\end{equation}
These terms are not redundant; they provide supervision at different levels.

\textbf{Outcome supervision: $\L_{\mathrm{diag}}$.}
Cross-entropy on the frozen diagnostic logits asks:
\emph{did the final evidence set answer correctly?}
It supervises the joint outcome but provides weak credit assignment among many candidate patches.

\textbf{Semantic supervision: $\L_{\mathrm{sem}}$.}
A frozen class-text prior scores how diagnostically discriminative a patch appears in the CONCH embedding space. (At inference the selector receives no class text, task identity, or label; during training, class-derived information enters only as supervision targets---this prior and $\L_{\rm diag}$. The attribution claim concerns the frozen non-selector parameters, not the absence of class-derived supervision.)
The selector distribution is softly aligned to this prior with KL divergence.
This is a weak semantic anchor, not the final diagnostic target.

\textbf{Counterfactual credit: $\L_{\mathrm{util}}$.}
For candidate patch $i$, define a marginal counterfactual gain
\begin{equation}
u_i
=
\L_{\mathrm{diag}}(E)
-
\L_{\mathrm{diag}}(E\cup\{x_i\}),
\label{eq:ui}
\end{equation}
where $E$ denotes the current evidence context.
Positive $u_i$ means that adding the candidate would reduce diagnostic loss; negative $u_i$ means it would be harmful.
The gains are treated as a no-gradient teacher distribution, and $\L_{\mathrm{util}}=\mathrm{KL}(\mathrm{softmax}(u)\,\Vert\,\mathrm{softmax}(s))$ trains the patch selector to rank candidates by measured usefulness.
Lowercase $u_i$ is therefore a \emph{candidate-level} teaching signal.

The value of the \emph{whole selected evidence set} is a different quantity:
\begin{equation}
U(\Psel)
=
\log C-\mathrm{CE}(\Psel,y).
\label{eq:bigU}
\end{equation}
It asks how much the complete selected set improves over an empty/uniform diagnostic state.
This distinction---candidate utility $u_i$ versus set utility $U$---becomes central in continual learning.

\begin{table}[t]
\centering
\footnotesize
\caption{\textbf{One study object, six training signals, matching questions.}
Within-task terms learn the current policy; continual terms preserve past policy behavior and function.}
\label{tab:questions}
\begin{tabular}{p{1.25cm}p{4.75cm}}
\toprule
Term & Question answered \\
\midrule
$\L_{\rm diag}$ & Did the selected set solve the current diagnosis? \\
$\L_{\rm sem}$ & Is the selector anchored to diagnostically discriminative semantics? \\
$\L_{\rm util}$ & Which candidate patch would actually help? \\
\midrule
$\L_{\rm KD}$ & Do I still score old evidence similarly? \\
$\L_{\rm eq}$ & If I select differently, is the new set at least as useful? \\
$\L_{\rm rep}$ & Can I still solve the old diagnostic problem? \\
\bottomrule
\end{tabular}
\end{table}

\subsection{One shared selector: LoRA consolidation}
\label{ssec:lora}

For task $\tau$, low-rank LoRA adapters are attached to the four linear layers of $\Fg$ and $\Fp$.
The base selector is frozen during the task and only the current adapters train.
At the task boundary, $W_{\tau}=W_{\tau-1}+B_\tau A_\tau$, which is an exact rewrite of the effective weight used immediately before merging.
The adapter is then reset for the next task.
Thus parameter count does not grow with $T$ and inference always uses one shared selector.

Crucially, \textbf{merging is not the anti-forgetting mechanism}.
The merge-only arm performs no better systematically than naive sequential fine-tuning.
LoRA supplies a clean, fixed-cost way to accumulate task updates; preservation requires additional constraints.

\subsection{Keeping looking well: preserve behavior and function}
\label{ssec:cl}

After a task is learned, PathSelect stores a bounded Selection Memory
$\Mem$ of at most 512 behavioral snapshots using reservoir sampling.
A snapshot stores slide/task identifiers, group scores $r_{\rm old}$, candidate indices, patch scores $s_{\rm old}$, and utilities; it stores \emph{no patch pixels and no patch features}.
Features are reloaded only when needed for training.

For each current sample, one old snapshot is replayed.
The continual objective is
\begin{align}
\L_{\mathrm{CL}}
={}&
\lambda_{\mathrm{kd}}
\left[
\mathrm{KL}(r_{\mathrm{old}}\Vert r_{\mathrm{new}})
+
\mathrm{KL}(s_{\mathrm{old}}\Vert s_{\mathrm{new}})
\right]
\nonumber\\
&+
\lambda_{\mathrm{eq}}
\max(0,U_{\mathrm{old}}-U_{\mathrm{new}})
+
\lambda_r \L_{\mathrm{diag}}^{\mathrm{rep}},
\label{eq:lcl}
\end{align}
with all three $\lambda$'s set to 1.

Equation~\ref{eq:lcl} deliberately separates three preservation questions.

\textbf{Behavioral continuity: $\L_{\rm KD}$.}
Dual-level distillation preserves the old \emph{soft selection policy}: group preferences and patch-score distributions.
It does not require the final hard top-$B$ set to be identical.

\textbf{Functional continuity: $\L_{\rm eq}$.}
Selection identity should be allowed to change if an alternative evidence set is equally or more useful.
The hinge
\[
\L_{\rm eq}=\max(0,U_{\rm old}-U_{\rm new})
\]
therefore acts as a \emph{floor, not a ceiling}: regression is penalized, improvement is free.
This is the central difference between ``remember the same evidence'' and ``retain the same evidence-selection competence.''

\textbf{Task competence: replay.}
$\L_{\mathrm{diag}}^{\mathrm{rep}}$ is simply the ordinary diagnostic loss on a replayed old slide.
Replay is a data mechanism rather than a new mathematical objective.
All within-task and continual terms update the \emph{same current selector LoRA} in one backward pass:
\[
\L_{\mathrm{total}}=\L_{\mathrm{ev}}+\L_{\mathrm{CL}}.
\]

% -----------------------------------------------------------------------------
\section{Experiments}
\label{sec:exp}

\subsection{Protocol and measurements}

We use four sequential TCGA organ tasks, ESCA$\rightarrow$RCC$\rightarrow$BRCA$\rightarrow$LUNG---the four-organ sequence used in prior continual WSI work \cite{huang2023conslide,gou2025qpmil}, traversed in reverse order---each a binary subtype task, for a final eight-class label space.
The train set contains 2,273 slides; test sizes are 15/76/93/95.
All runs use 5 epochs per task, learning rate $10^{-3}$, $B=8$, LoRA rank 4, $\beta_s{=}\beta_u{=}0.1$, unit $\lambda$'s, and seeds 0--4. Every comparison was registered with a two-sided decision rule before the corresponding runs: win counts over the five seeds decide, the incumbent design is retained on a split, and one registered expectation that the data refuted is reported as measured (Sec.~\ref{ssec:protocols}).

The central baselines are sequential full fine-tuning (SeqFT, A1), LoRA with boundary merging only (A2), replay (A3), replay+distillation (A4), and full PathSelect (A5=A4+utility hinge).
Per-task specialists and joint offline training are references rather than deployable continual baselines.

We measure outcome and selection separately.
Class-IL accuracy evaluates over all seen classes; task-IL accuracy restricts prediction to the task's own pair; leakage is the fraction predicted into another task's class set.
Selection Jaccard, $|\Psel_{\rm own}\cap\Psel_{\rm final}|/|\Psel_{\rm own}\cup\Psel_{\rm final}|$, measures hard identity retention between the set selected immediately after learning a task and the set selected at the end.
Jaccard is \emph{evaluation only}; it is not the distillation loss.
Functional retention is measured by
\[
\Delta U=U_{\rm final}-U_{\rm own}.
\]
With five seeds, comparisons are paired by seed and accompanied by win counts; 5/5 is called systematic, 4/5 directional, and no $p$-values are reported at this seed count.

\subsection{First establish the failure: forgetting where to look}

Before studying CL, learned selection must itself be meaningful.
In a preliminary budget sweep with a learned flat selector (inference only; seven budgets $\times$ four tasks), learned selection beat random selection in all 28 cells, and accuracy peaked at $B=8$ ($0.8797$ vs.\ $0.6607$ for random), supporting the small-evidence operating point.

The continual failure is then visible on three axes.
After unprotected sequential training, old-task class-IL accuracy collapses, the selected patch sets become nearly unrelated to their original evidence, and set utility can turn negative.
At the final checkpoint SeqFT has Selection Jaccard $0.0023$, near zero, and final class-IL accuracy $0.4774$.
This is not merely ``the classifier forgot'': the diagnostic pathway is frozen.
The learned observation policy itself has drifted.

\begin{table*}[t]
\centering
\footnotesize
\caption{\textbf{Continual-learning main comparison} (mean$\pm$sd over 5 seeds; current repository values).
LoRA merging provides a fixed-cost substrate but no systematic preservation.
Replay is strong; PathSelect adds behavioral and functional preservation and achieves the best final outcome among continual arms.}
\label{tab:main}
\begin{tabular}{lcccccc}
\toprule
Method &
Class-IL$\uparrow$ &
Forgetting$\downarrow$ &
Plasticity$\uparrow$ &
Sel. Jaccard$\uparrow$ &
Task-IL$\uparrow$ &
Leakage$\downarrow$ \\
\midrule
Per-task specialist (R1) & $0.8777{\pm}.020$ & --- & $0.8777$ & --- & $0.9027{\pm}.019$ & --- \\
Joint offline (R2)       & $0.7789{\pm}.028$ & --- & $0.7789$ & --- & $0.8498{\pm}.029$ & --- \\
\midrule
SeqFT (A1)               & $0.4774{\pm}.103$ & $0.5193$ & $0.8669$ & $0.0023$ & $0.8086{\pm}.061$ & $0.4408$ \\
LoRA-Merge (A2)          & $0.4466{\pm}.109$ & $0.5798$ & $0.8814$ & $0.0010$ & $0.8105{\pm}.036$ & $0.4756$ \\
$+$ Replay (A3)          & $0.7778{\pm}.015$ & $0.1102$ & $0.8514$ & $0.0864$ & $0.9073{\pm}.018$ & --- \\
$+$ Distillation (A4)    & $0.7972{\pm}.046$ & $0.1063$ & $0.8689$ & $\mathbf{0.1752}$ & $0.8969{\pm}.019$ & --- \\
\textbf{PathSelect (A5)} & $\mathbf{0.8239{\pm}.028}$ & $\mathbf{0.0539}$ & $0.8515$ & $0.1613$ & $\mathbf{0.9147{\pm}.014}$ & $\mathbf{0.1005}$ \\
\bottomrule
\end{tabular}
\end{table*}

\subsection{Main result: preservation, not merging, prevents forgetting}

Table~\ref{tab:main} gives the lifecycle of the method.
First, LoRA merging alone is not a CL solution: A2 does not systematically improve over SeqFT.
Second, replay is a strong baseline, raising final class-IL accuracy from $0.4774$ to $0.7778$.
Third, the complete preservation design reaches $0.8239$, a $+34.6$ point gain over SeqFT (5/5 seeds) and a $+4.61$ point gain over replay (5/5 under the main task order; this increment is order-sensitive and is reported as such below).
Forgetting falls from $0.5193$ (SeqFT) and $0.1102$ (replay) to $0.0539$.

Selection measurements explain \emph{how}.
Unprotected arms preserve almost none of the original evidence ($0.0023$/$0.0010$ Jaccard).
Replay recovers some identity ($0.0864$), and distillation-based methods recover an order of magnitude more ($0.1752$/$0.1613$).
Yet A4 has slightly higher Jaccard than A5 while A5 has better accuracy and forgetting.
Therefore exact identity retention is not the objective itself: \emph{remembering where} and \emph{retaining useful evidence} are different.

Utility retention reinforces this point.
The current experimental record reports $\Delta U=-16.8$ for the full method, versus approximately $-220$/$-301$ for unprotected arms.
Every method loses some old-task utility, but protected selectors lose an order of magnitude less.

\subsection{Do the preservation terms really do different jobs?}

\begin{table}[t]
\centering
\footnotesize
\caption{\textbf{Preservation-component ablation} (flat substrate, 5-seed means).
The terms are complementary rather than interchangeable.}
\label{tab:ablation}
\begin{tabular}{ccc|ccc}
\toprule
Replay & Distill. & Utility & Class-IL$\uparrow$ & Task-IL$\uparrow$ & Leak.$\downarrow$\\
\midrule
 & & & .447 & .811 & .476\\
 & \checkmark & & .600 & .853 & .323\\
 & & \checkmark & .809 & .885 & .091\\
\checkmark & & & .778 & .907 & .104\\
\checkmark & \checkmark & & .797 & .897 & .117\\
\checkmark & \checkmark & \checkmark & \textbf{.824} & \textbf{.915} & \textbf{.101}\\
\bottomrule
\end{tabular}
\end{table}

Table~\ref{tab:ablation} is best read mechanistically, not as a leaderboard.
Distillation protects \emph{how} the selector scores old evidence, yet distillation alone still suffers high cross-task leakage: preserving the selection policy does not by itself guarantee correct task attribution.
Replay restores old diagnostic supervision.
The utility hinge adds functional freedom: the selector need not reproduce the old hard evidence set if it can find another set with at least the same diagnostic value; its class-IL increment alone is order-sensitive, which we report rather than smooth over.
No single loss, and no pair, reproduces the full behavior.

The hierarchy strengthens the same preservation contrast rather than simply raising absolute accuracy.
In the flat selector, PathSelect-over-replay is only $+0.74{\pm}1.93$ task-IL (3/5).
Under hierarchical selection the contrast becomes $+3.28{\pm}2.40$ task-IL and $+5.76{\pm}3.42$ class-IL (both 5/5).
Removing group-level KD under the hierarchy costs $3.95$ task-IL and $3.71$ class-IL points (both 5/5), showing that group-level quota preservation and patch-level evidence preservation act at different resolutions.
We do not claim hierarchy raises the absolute ceiling: hierarchical A5 remains within noise of flat A5 and has somewhat higher leakage.
The pre-registered hierarchy gate asked whether preservation remained systematic when selection is decomposed into quota allocation and within-group choice; it did.

Finally, preservation design can matter more than buffer size within the tested range.
On the hierarchical configuration, PathSelect with 128 memory entries outperforms replay-only with 512 entries by $+3.06$ task-IL points (5/5); we report this memory-efficiency result for task-IL and this range only.
It supports the view that a memory of \emph{how the selector behaved} can be more efficient than replaying more samples.


\subsection{Are the update-protocol choices necessary?}
\label{ssec:protocols}

\begin{table}[t]
\centering\footnotesize
\caption{\textbf{Update-protocol ablations}: paired per-seed differences, variant minus reference (win counts count seeds favoring the variant; 0/5 means the reference wins every seed). Each registered probe was decided by a rule fixed before the run.}
\label{tab:abl}
\begin{tabular}{lccc}
\toprule
Variant $-$ reference & $\Delta$class-IL & $\Delta$task-IL & $\Delta$Leak. \\
\midrule
Warm start $-$ full & $-1.80$ (2/5) & $-0.96$ (1/5) & $+1.71$ (1/5) \\
Single adapter $-$ full & $-3.54$ (1/5) & $-2.06$ (1/5) & $+2.49$ (1/5) \\
Composition $-$ merge-only & $+24.6$ (5/5) & $+6.59$ (5/5) & $-24.9$ (5/5) \\
Composition $-$ full & $-13.1$ (0/5) & $-3.83$ (1/5) & $+12.6$ (0/5) \\
Damped merge $-$ full & $-0.92$ (3/5) & $-0.98$ (1/5) & $+1.15$ (1/5) \\
Single adapter $-$ full (hier) & $+0.79$ (4/5) & $-1.25$ (0/5) & $-3.25$ (5/5) \\
\bottomrule
\end{tabular}
\end{table}

Table~\ref{tab:abl} tests the protocol itself.
\textbf{Warm start} (full-parameter training of task 1 before switching to adapters) shows no systematic benefit, so the parameter-efficient incumbent stands.
\textbf{A single adapter that is never reset} loses on both accuracy axes on the flat substrate---fresh adapters supply capacity that grows with the task count---while on the hierarchical substrate the same paired comparison becomes a trade-off (systematic task-IL edge for fresh adapters, systematically lower leakage for the single adapter); we retain fresh adapters and report the interaction.
\textbf{Post-hoc composition} refuted our registered expectation: independently trained residuals, summed once, beat the bare sequential substrate on every metric in every seed, a result consistent with strong interference in unprotected sequential updates: independently learned residuals avoid cross-task overwriting (cf.\ task arithmetic \cite{ilharco2023task}). Composition is not the endpoint---its independence excludes replay, distillation targets require an accumulated model, and the full method beats it decisively (Table~\ref{tab:abl})---so the refined claim is that sequential consolidation earns its place only together with behavioral preservation, the component that requires sequence to define. Sum and average show no systematic difference under the registered comparison.
\textbf{Damped merging} ($\alpha{=}0.5$) helps nothing systematically; full merging stands.
Three of these interventions---warm start, damped merging, composition---raise selection Jaccard without improving diagnosis (composition attains the highest Jaccard of all arms, $0.275$, yet trails the full method by 13 points): selection identity is not selection usefulness, which is precisely the gap the utility hinge exists to close.
Three further components---a task-conditioned query, a stateful multi-round selector, and a group-level semantic prior---failed their registered gates and were removed from the method; the query variant retained a secondary leakage finding that we report in the extended version.

% -----------------------------------------------------------------------------
\section{Discussion and Limitations}
\label{sec:discussion}

PathSelect suggests that continual learning in budgeted perception should distinguish two failure modes that accuracy alone conflates.
A policy may change \emph{where} it looks, and the new evidence may or may not remain \emph{functionally equivalent}.
This is why Jaccard, KD, utility, and the utility hinge should not be collapsed into one notion.
KD preserves a soft behavior distribution; Jaccard measures hard identity after the fact.
Candidate utility $u_i$ teaches which patch is useful now; set utility $U$ evaluates whether the selected evidence remains useful across time.
The utility-preservation hinge uses $U$ as a lower bound, deliberately allowing better alternative evidence.

The study also reveals what is \emph{not} sufficient.
LoRA merging alone does not protect the selector.
Warm start, damped merging, and post-hoc composition each increase selection Jaccard without improving diagnosis (Sec.~\ref{ssec:protocols}), so ``choose the same patches'' is not a complete CL objective: selection stability is not sufficient for diagnostic preservation, and identity and function must be measured separately.
This motivates behavioral and functional preservation as complementary constraints rather than redundant regularizers.

Our claims remain bounded.
The benchmark uses four organ-separated TCGA tasks, one frozen foundation model (CONCH), one main budget ($B=8$), known task boundaries during training, and five seeds.
The tasks are highly separable (a linear probe separates them at 98.21\%), and explicit task conditioning failed its registered gate, so harder within-organ or mixed-task streams may yield different conclusions.
PathSelect is also a one-shot evidence-acquisition policy, not a multi-step navigation agent.
Its task-free shared selector can serve as a foundation for continual evidence-navigation systems, but sequential stateful acquisition is left to future work.

% -----------------------------------------------------------------------------
\section{Conclusion}

PathSelect reframes continual WSI learning around a simple but under-measured question: after learning new diagnostic tasks, does the model still know \emph{where to look}?
By freezing the diagnostic pathway, we isolate the evidence selector as the study object.
Within tasks, semantic and counterfactual signals teach the selector to find useful evidence.
Across tasks, a single merged selector is preserved at two complementary levels: distillation maintains selection behavior, while a utility-preservation hinge allows evidence identity to change without allowing usefulness to regress; replay maintains old-task competence.
The resulting method substantially reduces forgetting at zero inference-time memory or routing cost.
More broadly, the results argue that continual learning for budgeted perception should preserve not only what a model predicts, but also the policy by which it acquires the evidence used to predict.

\bibliographystyle{IEEEbib}
\bibliography{refs}

\end{document}
